\section{Study Context and Sampling}

Vaccination is widely recognized as a cornerstone of infectious-disease prevention, yet real-world vaccination behavior is not simply a binary choice between uptake and refusal. A central policy challenge is heterogeneity in timing: many individuals delay vaccination even when vaccines are available, thereby remaining exposed to infection risk for longer periods and weakening the effective coverage of immunization programs \citep{Halilova2022, Halilova2023}.

Empirically, this dissertation studies these timing trade-offs in a Chinese city---Wuhan---during the broader COVID-era immunization environment, where vaccination decisions were highly salient and program logistics could generate meaningful variation in waiting time. This setting provides a salient context for measuring how individuals value timeliness when delayed protection implies a longer exposure window. Because delayed protection is evaluated in a high-salience context, the estimated valuation of waiting provides an informative benchmark for other immunization settings, where perceived risk and salience may differ.

From a policy perspective, these delays matter because they shift protection into the future while infection risk persists in the present, slowing the epidemiological and social benefits of immunization. Consequently, vaccination timing can be understood as an intertemporal health investment decision in which immediate costs (time, inconvenience, perceived side effects) are weighed against delayed and uncertain health benefits \citep{Grossman1972, Meissner2025, Strickland2022}. This intertemporal framing is especially relevant in infectious-disease settings where ``waiting'' has biological meaning: delaying vaccination extends the period of vulnerability and can amplify perceived uncertainty and behavioral burdens.

At the same time, behavioral research emphasizes that vaccination decisions are shaped by psychological and contextual determinants---such as confidence, perceived risks, and perceived barriers---which can interact with intertemporal trade-offs when protection is delayed \citep{Halilova2024, Grytten2024}. These insights imply that policies aimed at improving vaccination timeliness may require more than information provision or supply expansion alone; they may also involve reducing waiting time, offering compensation or incentives, and targeting high-burden groups whose delay sensitivity is systematically greater.

Empirically, identifying the valuation of delayed protection is challenging using observational timing data alone because realized waiting time is jointly determined by preferences, constraints, and program rules. To isolate preferences for timeliness in a policy-relevant way, this dissertation uses a discrete choice experiment (DCE) that embeds waiting time to protection directly as an attribute of vaccination alternatives, allowing estimation of trade-offs between timeliness and other vaccine attributes under controlled counterfactual scenarios \citep{Louviere2000, Train2009, DCEreview2025}. This approach supports two complementary policy outputs: structural estimates of how individuals discount delayed protection and monetary translations of delay aversion that can be used to compare ``speeding up access'' versus ``compensation'' strategies.
